\chapter{Distance, error, and the meaning of a limit}\label{ch:limits}
\goal{Describe distance by three simple rules, so that one argument can work for numbers, vectors, and other objects. Replace the impression that a sequence ``gets close'' to a point by a precise promise that covers every positive tolerance. See why steps that keep getting smaller need not make a sequence settle down, while steps that shrink by a fixed factor do. Recognize completeness as the property that guarantees a limit inside the space for every sequence whose terms eventually all lie as close together as we like.}

\section{Distance begins with absolute value}\label{ch11:sec:distance}
Recall from Section~\ref{sec:magnitude} that the absolute value $|x|$ of a real number $x$ is $x$ when $x\ge0$ and $-x$ when $x<0$. For instance, $|4|=4$ and $|-3|=-(-3)=3$. The \emph{distance} between two real numbers $x$ and $y$ is $|x-y|$, the length of the stretch of the number line between them. For $x=-2$ and $y=3$ it is $|-2-3|=|-5|=5$: walking from $-2$ to $3$ takes two units to reach $0$ and three more to reach $3$. The order of the two numbers does not matter, since $|y-x|=|3-(-2)|=|5|=5$ as well. In general $y-x=-(x-y)$, and changing the sign of a number does not change its absolute value, so $|y-x|=|x-y|$.

This chapter starts from the familiar distance $|x-y|$ and asks which of its properties our arguments actually use. Once those few properties are written down, any argument that relies only on them works just as well for other objects, such as the vectors of Chapter~\ref{ch:linear}, provided they come with a distance that has the same properties.

The most important of these properties is the \emph{triangle inequality}: for all real numbers $x$ and $y$,
\[
|x+y|\le|x|+|y|.
\]
For example, $x=5$ and $y=-2$ give $|x+y|=3$, which is less than $|x|+|y|=7$. When $x$ and $y$ are both nonnegative, or both nonpositive, the two sides are equal; for instance, $x=-5$ and $y=-2$ give $|x+y|=|-7|=7$ and $|x|+|y|=5+2=7$.

To prove the inequality, recall another rule from Section~\ref{sec:magnitude}: for a number $b\ge0$, the statement $|x|\le b$ holds exactly when $-b\le x\le b$. Taking $b=|x|$, the true statement $|x|\le|x|$ becomes $-|x|\le x\le|x|$. In the same way, $-|y|\le y\le|y|$. Two inequalities in the same direction may be added: if $p\le r$ and $s\le t$, then $p+s\le r+s\le r+t$. Adding $-|x|\le x$ and $-|y|\le y$ gives $-|x|-|y|\le x+y$, and adding $x\le|x|$ and $y\le|y|$ gives $x+y\le|x|+|y|$. Together,
\[
-(|x|+|y|)\le x+y\le|x|+|y|.
\]
Now use the same rule in the other direction, with $x+y$ in place of $x$ and the nonnegative number $b=|x|+|y|$. It turns this pair of inequalities into $|x+y|\le|x|+|y|$.

The triangle inequality is most often used in a second form. Given three real numbers $a$, $b$ and $c$, take $x=a-b$ and $y=b-c$. Then $x+y=a-c$, because the two copies of $b$ cancel, and the triangle inequality becomes
\[
|a-c|\le|a-b|+|b-c|.
\]
Read as a statement about distances, it says that the distance from $a$ to $c$ is at most the distance from $a$ to $b$ plus the distance from $b$ to $c$. Going from $a$ to $c$ by way of a stopping point $b$ is never shorter than going directly. For $a=0$, $b=5$ and $c=2$, the direct distance is $2$, while the route through $b$ has length $5+3=8$. When $b$ lies between $a$ and $c$, as in $a=0$, $b=1$, $c=2$, the two sides are equal.

This second form explains why the triangle inequality appears in almost every error estimate in this chapter. To show that two quantities $a$ and $c$ are close, it is often easier to find a third quantity $b$ that is close to each of them, and then to add the two small distances.

Here is a first error estimate. Suppose $x$ approximates a number $a$ with error at most $\eps$, and $y$ approximates a number $b$ with error at most $\delta$, where $\delta$ (the Greek letter delta) is a second tolerance: $|x-a|\le\eps$ and $|y-b|\le\delta$. How well does $x+y$ approximate $a+b$? Regroup the difference as $(x+y)-(a+b)=(x-a)+(y-b)$ and apply the triangle inequality:
\[
|(x+y)-(a+b)|\le |x-a|+|y-b|\le\eps+\delta.
\]
So when approximations are added, their error bounds add. For example, take $a=2$, $b=5$, $\eps=0.01$ and $\delta=0.02$. Then $x+y$ is within $0.03$ of $7$. This bound is reached when both errors point the same way: $x=2.01$ and $y=5.02$ give $x+y=7.03$. When the errors point in opposite directions, they partly cancel: $x=2.01$ and $y=4.98$ give $x+y=6.99$, only $0.01$ away from $7$. Since we know only the sizes of the two errors, not their directions, $\eps+\delta$ is the best guarantee available: it promises that the error is never larger, not that it is always this large.

\pauseq{Suppose $|x-a|\le\eps$ and $|y-b|\le\delta$. Does it follow that $|(x-y)-(a-b)|\le\eps-\delta$? If not, what bound is correct?}{No. Take $a=b=0$, $x=\eps$ and $y=-\delta$. Both hypotheses hold, and $(x-y)-(a-b)=\eps+\delta$, which is larger than $\eps-\delta$ because $\delta>0$. The correct bound is $\eps+\delta$. To prove it, write $(x-y)-(a-b)=(x-a)+\bigl(-(y-b)\bigr)$ and apply the triangle inequality. Since $|-(y-b)|=|y-b|$, the result is $|(x-y)-(a-b)|\le|x-a|+|y-b|\le\eps+\delta$. Subtracting two approximations does not subtract their errors; in the worst case the errors still add.}

\section{A metric keeps only the properties we use}\label{ch11:sec:metric}
Let $X$ be a set, whose members we call \emph{points}, even when they are numbers or vectors. A distance on $X$ should be a function $d$ from $X\times X$ to $\RR$ (Sections~\ref{ch03:sec:sets} and~\ref{ch03:sec:functions}).
\par\noindent\begin{minipage}{\linewidth}
\begin{definition}[Metric]\label{def:metric}
A \emph{metric} on a set $X$ assigns a real number $d(x,y)$ to every pair of points $x,y\in X$ and satisfies these requirements for all $x,y,z\in X$:
\begin{enumerate}
\item $d(x,y)\ge0$, and $d(x,y)=0$ holds exactly when $x=y$.
\item $d(x,y)=d(y,x)$.
\item $d(x,z)\le d(x,y)+d(y,z)$.
\end{enumerate}
A set together with a chosen metric is called a \emph{metric space} and is written $(X,d)$. The number $d(x,y)$ is called the \emph{distance} between $x$ and $y$.
\end{definition}
\end{minipage}\par

The first requirement says, in particular, that two different points are a positive distance apart. The second is called \emph{symmetry}, and the third is the triangle inequality of Section~\ref{ch11:sec:distance}, now demanded of $d$. To show that a new formula is a metric, we must \emph{recognize} it as one in the sense of Section~\ref{ch01:sec:definition}, by checking all three requirements; calling a formula a distance does not make it one.

\pauseq{On the real line, let $d(x,y)=(x-y)^2$. Is $d$ a metric?}{No. The first two requirements hold: a square is never negative, $(x-y)^2=0$ exactly when $x=y$, and $(y-x)^2=(x-y)^2$. The triangle inequality fails. Take $x=0$, $y=1$ and $z=2$. Then $d(x,z)=4$, while $d(x,y)+d(y,z)=1+1=2$, so the route through $y$ comes out \emph{shorter} than the direct route. One counterexample is enough.}

\subsection*{Examples}
The real line with $d(x,y)=|x-y|$ is a metric space. The absolute value is never negative, and $|x-y|=0$ exactly when $x=y$; symmetry and the triangle inequality were proved in Section~\ref{ch11:sec:distance}. When we treat the real line as a metric space without naming a metric, this is the metric we mean.

The set $\RR^n$ of vectors $x=(x_1,\ldots,x_n)$ from Section~\ref{ch10:sec:vectors} carries several useful metrics. The first is the \emph{Euclidean distance} $d_2(x,y)=\norm{x-y}_2$, the length of the difference vector; in the plane it is the straight-line distance, so from $(0,0)$ to $(3,4)$ it is $\sqrt{3^2+4^2}=5$. It is a metric: a length is never negative and is zero only for the zero vector (Section~\ref{ch10:sec:inner}), which $x-y$ is exactly when $x=y$; changing the sign of every coordinate of $x-y$ does not change the squares under the root, so $d_2(y,x)=d_2(x,y)$; and the triangle inequality $\norm{x-z}_2\le\norm{x-y}_2+\norm{y-z}_2$ is proved at the end of Section~\ref{ch10:sec:cs}.

Another useful choice is
\[
d_\infty(x,y)=\max_{1\le j\le n}|x_j-y_j|,
\]
the largest of the $n$ numbers $|x_1-y_1|,\ldots,|x_n-y_n|$; the subscript $\infty$ is just part of the traditional name. This metric records the worst error in any single coordinate. For $x=(1,5)$ and $y=(4,3)$, the coordinate differences have sizes $3$ and $2$, so $d_\infty(x,y)=3$, while $d_2(x,y)=\sqrt{13}\approx3.61$.

To prove the triangle inequality for $d_\infty$, let $x,y,z\in\RR^n$ and fix a coordinate $j$. The number $|x_j-y_j|$ is one of those whose largest is $d_\infty(x,y)$, so $|x_j-y_j|\le d_\infty(x,y)$; likewise $|y_j-z_j|\le d_\infty(y,z)$. With the triangle inequality for real numbers,
\[
|x_j-z_j|\le|x_j-y_j|+|y_j-z_j|\le d_\infty(x,y)+d_\infty(y,z).
\]
The right side does not depend on $j$, so it is at least each of the numbers $|x_j-z_j|$, and hence at least their largest: $d_\infty(x,z)\le d_\infty(x,y)+d_\infty(y,z)$. Exercise~\ref{ex:11.3} checks the other two requirements.

Which metric to use depends on the question. If an instrument must get \emph{every} coordinate right to within $0.01$, the condition is $d_\infty(x,y)\le0.01$; a bound on $d_2$ instead controls all the coordinate errors together, through the sum of their squares. The reading $x=(0.01,0.01)$ of the true point $y=(0,0)$ meets the first condition, since $d_\infty(x,y)=0.01$, but fails $d_2(x,y)\le0.01$, since $d_2(x,y)\approx0.0141$.

\subsection*{Norms, subsets, and spaces without addition}
Both metrics on $\RR^n$ measure the size of a difference vector: $d_2(x,y)=\norm{x-y}_2$, and $d_\infty(x,y)=\norm{x-y}_\infty$ with $\norm{x}_\infty=\max_j|x_j|$. Such measures of size are called norms. Precisely, a \emph{norm} on $\RR^n$ assigns to each vector $x$ a real number $\norm{x}$ such that, for all vectors $x,y$ and every real number $c$,
\begin{enumerate}
\item $\norm{x}\ge0$, with equality exactly when $x$ is the zero vector;
\item $\norm{cx}=|c|\,\norm{x}$;
\item $\norm{x+y}\le\norm{x}+\norm{y}$.
\end{enumerate}
The Euclidean length is a norm, by Sections~\ref{ch10:sec:inner} and~\ref{ch10:sec:cs}. So is $\norm{x}_\infty$: it is never negative and is $0$ exactly when every coordinate is $0$, multiplying every coordinate by $c$ multiplies the largest coordinate size by $|c|$, and the third property is proved one coordinate at a time, like the triangle inequality for $d_\infty$. Conversely, every norm gives a metric $d(x,y)=\norm{x-y}$: the first requirement comes from the first property of the norm, since $x-y$ is the zero vector exactly when $x=y$; symmetry from the second with $c=-1$, since $y-x=(-1)(x-y)$; and the triangle inequality from the third, since $x-z=(x-y)+(y-z)$.

A metric, however, needs only one number $d(x,y)$ for each pair of points, not a way to add points. So the rest of the chapter works in a general metric space, and its arguments apply also to sets in which addition makes no sense, or can lead outside the set. The most common such set is part of a larger metric space. If $(X,d)$ is a metric space and $C\subseteq X$, then $d$, used only for pairs of points of $C$ (the restriction of $d$ to $C\times C$, in the language of Section~\ref{ch03:sec:injsurj}), is a metric on $C$: the three requirements hold for all points of $X$, so in particular for all points of $C$. For example, the interval $[0,1]$ with distance $|x-y|$ is a metric space, even though a sum of its points can leave it, as $0.7+0.6=1.3$ does.

\section{Convergence promises eventual control}\label{ch11:sec:convergence}
A \emph{sequence} in a metric space $(X,d)$ is an endless list $x_0,x_1,x_2,\ldots$ of points of $X$, one for each natural number, as in Section~\ref{ch05:sec:recurrence} for $X=\RR$. Here the subscript counts terms: $x_n$ is a whole point of $X$, not a coordinate of a vector. The terms $x_N,x_{N+1},x_{N+2},\ldots$ from some index $N$ onward form a \emph{tail} of the sequence.
\par\noindent\begin{minipage}{\linewidth}
\begin{definition}[Convergence]\label{ch11:def:convergence}
A sequence $x_0,x_1,x_2,\ldots$ in a metric space $(X,d)$ \emph{converges} to a point $x_*\in X$ if, for every tolerance $\eps>0$, there is a threshold $N$ such that every term from $x_N$ onward lies at distance less than $\eps$ from $x_*$:
\[
\forall\eps>0\ \exists N\in\NN\ \forall n\ge N:\quad d(x_n,x_*)<\eps.
\]
The point $x_*$ (the star is part of the name) is then called a \emph{limit} of the sequence, and we write $x_n\to x_*$, read ``$x_n$ tends to $x_*$.'' A sequence that converges to some point of $X$ is called \emph{convergent}.
\end{definition}
\end{minipage}\par

Read the quantifiers as the skeptic--prover game of Section~\ref{ch02:sec:quantifiers}, with $x_*$ fixed in advance. The skeptic names a tolerance $\eps$; the prover, knowing $\eps$, names a threshold $N$; the skeptic then picks \emph{any} index $n\ge N$, and the prover wins if $d(x_n,x_*)<\eps$. The sequence converges to $x_*$ when the prover can always win. A smaller tolerance may force a larger threshold, which is allowed because $N$ is chosen after $\eps$; but a single $N$ must handle every $n\ge N$ at once.

\subsection*{Two sequences that converge}
The statement $x_n\to x_*$ does not say that the sequence ever reaches its limit. Take $x_n=1/(n+1)$ on the real line, so that the sequence is $1,\tfrac12,\tfrac13,\ldots$; its terms are all positive, yet they tend to $0$. Let $\eps>0$ be given. By the Archimedean property (Section~\ref{ch02:sec:quantifiers}) we may choose an integer $N>1/\eps$; it is positive because $1/\eps>0$. Now let $n\ge N$. Then $n+1>N>0$, and taking reciprocals reverses the order, so $1/(n+1)<1/N$. Taking reciprocals in $N>1/\eps$ gives $1/N<\eps$. Together,
\[
d(x_n,0)=\left|\frac1{n+1}-0\right|=\frac1{n+1}<\frac1N<\eps.
\]
This proves $x_n\to0$. For $\eps=0.01$ the recipe asks for an integer $N>100$, such as $N=101$; the threshold $N=100$ would also work, since $x_{100}=1/101<0.01$. The definition asks only for \emph{some} threshold that works, not the smallest one.

Next, the sequence $x_n=1-2^{-n}$ of Section~\ref{ch05:sec:recurrence}, described there as coming as close to $1$ as we like, does converge to $1$. The distance from $x_n$ to $1$ is $2^{-n}$. To compare it with $1/(n+1)$, we show by induction that $2^n\ge n+1$ for every $n\in\NN$. At $n=0$ both sides equal $1$. If $2^n\ge n+1$ for some $n$, then
\[
2^{n+1}=2\cdot2^n\ge2(n+1)=2n+2\ge n+2,
\]
which is the claim for $n+1$. Taking reciprocals gives $2^{-n}\le1/(n+1)$. So if $\eps>0$ and $N$ is an integer with $N>1/\eps$, every $n\ge N$ satisfies $|x_n-1|\le1/(n+1)<\eps$. This proves $x_n\to1$.

Because the definition speaks only about tails, changing finitely many terms of a sequence affects neither whether it converges nor its limit. If the first thousand terms of $1/(n+1)$ were replaced by arbitrary numbers, a threshold that is at least $1000$ and larger than $1/\eps$ would still work.

\subsection*{What failure to converge looks like}
Negating the definition with the rules of Section~\ref{ch02:sec:negation} shows what it means for a sequence \emph{not} to converge to $x_*$:
\[
\exists\eps>0\ \forall N\in\NN\ \exists n\ge N:\quad d(x_n,x_*)\ge\eps.
\]
There is one fixed tolerance such that, however late we start, some later term is at least that far from $x_*$.

Consider the sequence $0,1,0,1,\ldots$, with $x_n=0$ for even $n$ and $x_n=1$ for odd $n$. It returns to $0$ again and again, but convergence needs a whole tail near the limit. Take $\eps=1/2$. For every threshold $N$, the odd index $n=2N+1$ is at least $N$, and $d(x_n,0)=1\ge1/2$. So the sequence does not converge to $0$.

\pauseq{Does the sequence $0,1,0,1,\ldots$ converge to $1/2$? Does it converge to any real number at all?}{It does not converge to $1/2$: every term is at distance exactly $1/2$ from $1/2$, so for $\eps=1/2$ no term is close enough. In fact it converges to no real number $L$. Suppose it did, and take $\eps=1/2$. Some threshold $N$ would make every term from $x_N$ onward lie less than $1/2$ away from $L$, including $x_{2N}=0$ and $x_{2N+1}=1$. The triangle inequality through $L$ would then give $1=|1-0|\le|1-L|+|L-0|<\tfrac12+\tfrac12=1$, which is impossible.}

\section{Limits are unique}\label{ch11:sec:unique}
Can a sequence converge to two different points at once? The proof below uses only the three requirements of a metric, so it settles the question in every metric space.
\par\noindent\begin{minipage}{\linewidth}
\begin{proposition}\label{thm:limitunique}
A sequence in a metric space has at most one limit: if $x_n\to a$ and $x_n\to b$, then $a=b$.
\end{proposition}
\end{minipage}\par

\begin{proof}
Suppose $x_n\to a$ and $x_n\to b$, and assume, for contradiction, that $a\ne b$. By the first requirement of a metric, $D=d(a,b)$ is then positive, and so is the tolerance $D/3$. Convergence to $a$ supplies a threshold $N_1$ such that $d(x_n,a)<D/3$ for every $n\ge N_1$, and convergence to $b$ supplies a threshold $N_2$ such that $d(x_n,b)<D/3$ for every $n\ge N_2$. Take a single index $n$ at least as large as both thresholds, for instance the larger of $N_1$ and $N_2$. The triangle inequality with $x_n$ as the stopping point, together with the symmetry $d(a,x_n)=d(x_n,a)$, gives
\[
D=d(a,b)\le d(a,x_n)+d(x_n,b)<\frac D3+\frac D3=\frac{2D}3.
\]
So $D<2D/3$, that is, $D/3<0$, which contradicts $D>0$. Therefore $a=b$.
\end{proof}

Any tolerance $\eps>0$ with $2\eps\le D$ would do in place of $D/3$, since the proof then gives $D<2\eps\le D$. The key step is the choice of one index beyond both thresholds. The triangle inequality needs a single point $x_n$ close to both $a$ and $b$, as a stop on the route from $a$ to $b$; a term close to $a$ and a \emph{different} term close to $b$ would give two pieces that do not join up.

Because of the proposition, we may speak of \emph{the} limit of a convergent sequence: $1/(n+1)$, for example, converges to $0$ and to no other real number.

\section{Cauchy sequences do not need a known destination}\label{ch11:sec:cauchy}
To prove that a sequence converges, the definition asks us to name the limit first, and often we do not know it. In Chapter~\ref{ch:contraction}, for example, a sequence is produced by applying the same rule again and again, and the point it approaches is exactly the unknown we are trying to find. The following condition, named after the French mathematician Augustin-Louis Cauchy (1789--1857), speaks only about the terms themselves.
\par\noindent\begin{minipage}{\linewidth}
\begin{definition}[Cauchy sequence]\label{ch11:def:cauchy}
A sequence $x_0,x_1,x_2,\ldots$ in a metric space $(X,d)$ is a \emph{Cauchy sequence} if, for every tolerance $\eps>0$, there is a threshold $N$ such that any two terms from $x_N$ onward are less than $\eps$ apart:
\[
\forall\eps>0\ \exists N\in\NN\ \forall m,n\ge N:\quad d(x_m,x_n)<\eps.
\]
\end{definition}
\end{minipage}\par

For convergence, the late terms had to be close to one fixed point $x_*$; here they have to be close to one another, and no limit appears. In the game, the skeptic, after seeing $N$, may now pick \emph{any two} indices $m,n\ge N$.

\par\noindent\begin{minipage}{\linewidth}
\begin{proposition}\label{ch11:prop:convergent-cauchy}
Every convergent sequence in a metric space is a Cauchy sequence.
\end{proposition}
\end{minipage}\par

\begin{proof}
Suppose $x_n\to x_*$, and let $\eps>0$ be given. Apply the definition of convergence with the smaller tolerance $\eps/2$. It supplies a threshold $N$ such that $d(x_n,x_*)<\eps/2$ for every $n\ge N$. Now let $m,n\ge N$, so that both $x_m$ and $x_n$ are less than $\eps/2$ away from $x_*$. The triangle inequality with $x_*$ as the stopping point, together with symmetry, gives
\[
 d(x_m,x_n)\le d(x_m,x_*)+d(x_*,x_n)
 <\frac\eps2+\frac\eps2=\eps.
\]
So the threshold $N$ works for every pair of indices from $N$ on, and the sequence is Cauchy.
\end{proof}

We used the tolerance $\eps/2$ because the triangle inequality adds two distances. Splitting a tolerance into pieces in this way is a standard move that will return in later chapters.

The converse, ``every Cauchy sequence converges,'' is not true in every metric space: the terms crowd together, but the point they crowd around may be missing from the space (Section~\ref{ch11:sec:complete}).

\subsection*{Small steps are not enough}
A sequence whose consecutive terms get closer and closer need not be Cauchy. Consider the \emph{harmonic sums}
\[
s_0=0,\qquad s_n=1+\frac12+\frac13+\cdots+\frac1n=\sum_{j=1}^n\frac1j\quad(n\ge1).
\]
Each step adds one more term, so $s_{n+1}-s_n=1/(n+1)$, and Section~\ref{ch11:sec:convergence} showed that these step sizes tend to $0$. Yet the sums do not settle down. For $n\ge1$, the difference $s_{2n}-s_n=\frac1{n+1}+\cdots+\frac1{2n}$ has $n$ terms, each at least the smallest one, $1/(2n)$. For instance, $s_8-s_4=\frac15+\frac16+\frac17+\frac18\ge4\cdot\frac18=\frac12$. In general,
\[
s_{2n}-s_n=\sum_{j=n+1}^{2n}\frac1j\ge n\cdot\frac1{2n}=\frac12.
\]
However late we start, doubling the index adds at least $1/2$. To turn this into a proof, negate the definition of a Cauchy sequence, as we did for convergence:
\[
\exists\eps>0\ \forall N\in\NN\ \exists m,n\ge N:\quad d(x_m,x_n)\ge\eps.
\]
Take $\eps=1/2$, and let $N$ be any proposed threshold. Choose $n=\max\{N,1\}$; we need $n\ge1$ because the estimate $s_{2n}-s_n\ge\tfrac12$ holds only from $n=1$ on (for $n=0$ the difference is $0$). Then choose $m=2n$. Both indices are at least $N$, and $|s_m-s_n|=s_{2n}-s_n\ge1/2$. So no threshold works for the tolerance $1/2$, and the sequence is not Cauchy. By the contrapositive (Section~\ref{ch02:sec:contrapositive}) of Proposition~\ref{ch11:prop:convergent-cauchy}, the harmonic sums therefore do not converge.

\section{Completeness excludes missing limits}\label{ch11:sec:complete}
Section~\ref{ch11:sec:cauchy} showed that every convergent sequence is Cauchy and noted that the converse can fail. The spaces in which it never fails have a name.
\par\noindent\begin{minipage}{\linewidth}
\begin{definition}[Complete metric space]\label{ch11:def:complete}
A metric space $(X,d)$ is \emph{complete} if every Cauchy sequence in $X$ converges to a point of $X$.
\end{definition}
\end{minipage}\par

Informally, a complete space has no holes: whenever the terms of a sequence eventually all lie as close to one another as we like, there is a point of the space that they approach. Completeness is a property of a whole space, not of a single sequence. To use it, we prove that one particular sequence is Cauchy, and the definition hands us a limit in the space. To show that a space is \emph{not} complete, it is enough to find one Cauchy sequence in it that has no limit in the space.

\subsection*{An interval with a missing endpoint}
The interval $(0,1)$, with distance $|x-y|$, is not complete. Consider the sequence $x_n=2^{-(n+1)}$, that is, $\tfrac12,\tfrac14,\tfrac18,\ldots$. Every term lies in $(0,1)$, because $0<2^{-(n+1)}\le\tfrac12$.

First, the sequence is Cauchy. On the real line it converges to $0$: the inequality $2^k\ge k+1$ of Section~\ref{ch11:sec:convergence}, with $k=n+1$, gives $|x_n-0|\le1/(n+2)<1/(n+1)$, and the argument given there for $1/(n+1)$ shows that these distances fall below any $\eps>0$ from some threshold on. So the sequence is Cauchy on the real line, by Proposition~\ref{ch11:prop:convergent-cauchy}. The Cauchy condition involves only distances between terms, which are the same in $(0,1)$ as in $\RR$, so the sequence is Cauchy in $(0,1)$ as well.

Second, the sequence has no limit in $(0,1)$. A limit $L\in(0,1)$ would also be a limit on the real line, since distances in $(0,1)$ are distances in $\RR$. As the sequence already converges to $0$ there, Proposition~\ref{thm:limitunique} would give $L=0$, which is not in $(0,1)$. The only possible destination of the sequence has been removed from the space.

\subsection*{The real line}
The real line, with distance $|x-y|$, is complete. We accept this as the last of the four starting facts about the real numbers listed in \emph{To the reader}; a proof would require constructing the real numbers from simpler objects. The closed intervals proved complete below, and every later argument that needs a limit to exist without naming it in advance, rest on this one fact.

The rational numbers show what the assumption rules out. Cutting off the decimal expansion of $\sqrt2=1.41421\ldots$ after $k$ decimal places gives rational numbers $t_0=1$, $t_1=1.4$, $t_2=1.41$, $t_3=1.414$, and so on, which crowd around $\sqrt2$. Since $\sqrt2$ is not rational (Theorem~\ref{thm:sqrt2}), $\QQ$, with distance $|x-y|$, has a hole there: it is not complete.

To make this precise, first check that $t_k\to\sqrt2$ on the real line. The difference $\sqrt2-t_k$ consists of the digits of $\sqrt2$ after the $k$th decimal place (for instance, $\sqrt2-t_2=0.00421\ldots$). Its first $k$ decimal digits are zeros, so $0\le\sqrt2-t_k\le10^{-k}$. Since $10^k\ge2^k\ge k+1$, this is at most $1/(k+1)$, and the argument of Section~\ref{ch11:sec:convergence} for $1/(k+1)$ gives $t_k\to\sqrt2$. Now argue as for $(0,1)$. By Proposition~\ref{ch11:prop:convergent-cauchy}, the sequence is Cauchy on the real line, and so also in $\QQ$, since the distances $|t_m-t_n|$ are the same. If it converged in $\QQ$ to a rational number $L$, it would converge to $L$ on the real line too, and Proposition~\ref{thm:limitunique} would give $L=\sqrt2$, which is not rational.

\subsection*{Closed subsets of a complete space}
Completeness passes from a space to many of its subsets, each given the same distance as in Section~\ref{ch11:sec:metric}.
\par\noindent\begin{minipage}{\linewidth}
\begin{definition}[Closed subset]\label{ch11:def:closed}
A subset $C$ of a metric space $(X,d)$ is \emph{closed} in $X$ if, whenever a sequence of points of $C$ converges to a point of $X$, that limit lies in $C$.
\end{definition}
\end{minipage}\par

In words, you cannot leave $C$ by taking limits.
\par\noindent\begin{minipage}{\linewidth}
\begin{proposition}\label{ch11:prop:closed-complete}
If $(X,d)$ is complete and $C$ is closed in $X$, then $C$, with the same distance $d$, is complete.
\end{proposition}
\end{minipage}\par

\begin{proof}
Take a Cauchy sequence of points of $C$. The distances between its terms are the same in $C$ as in $X$, so it is also Cauchy in $X$. Since $X$ is complete, it converges to some point of $X$. Since $C$ is closed, that point lies in $C$, and the distances from the terms to it are again the same in $C$ as in $X$. So the sequence converges to a point of $C$.
\end{proof}

The interval $(0,1]$ shows what goes wrong without closedness. The points $1/(n+1)$ lie in $(0,1]$, but their limit $0$ on the real line does not, so $(0,1]$ is not closed in $\RR$; and exactly as for $(0,1)$, this sequence is Cauchy in $(0,1]$ with no limit there, so $(0,1]$ is not complete.

The closed interval $[0,1]$, on the other hand, is closed in $\RR$, and therefore complete. To prove that it is closed, suppose points $x_n\in[0,1]$ converge on the real line to a number $L$. We show that $L<0$ and $L>1$ are both impossible, using the rule from Section~\ref{sec:magnitude} that, for $\eps>0$, the inequality $|x_n-L|<\eps$ means $L-\eps<x_n<L+\eps$.

If $L<0$, take half the distance from $L$ to $0$ as the tolerance: $\eps=-L/2>0$. Convergence supplies a threshold $N$ such that $|x_n-L|<-L/2$ for every $n\ge N$, and in particular
\[
x_N<L+\left(-\frac L2\right)=\frac L2<0,
\]
contradicting $x_N\in[0,1]$. If $L>1$, the positive tolerance $\eps=(L-1)/2$ gives in the same way a term with $|x_N-L|<\eps$, so
\[
 x_N>L-\frac{L-1}2=\frac{L+1}2>1,
\]
because $L>1$ gives $L+1>2$. This again contradicts $x_N\in[0,1]$. So $L\in[0,1]$. Hence $[0,1]$ is closed in the complete real line, and by Proposition~\ref{ch11:prop:closed-complete} it is complete.

The same argument, with $a$ and $b$ in place of $0$ and $1$ and the tolerances $(a-L)/2$ and $(L-b)/2$, shows that every closed interval $[a,b]$ with $a<b$ is complete. For instance, if $L<a$, a late term satisfies $x_N<L+(a-L)/2=(L+a)/2<a$.

\pauseq{Is the set $[0,\infty)$ of nonnegative real numbers, with distance $|x-y|$, complete?}{Yes. If nonnegative numbers converged on the real line to some $L<0$, the first case of the argument for $[0,1]$ would produce a negative term. So $L\ge0$, the set $[0,\infty)$ is closed in $\RR$, and by Proposition~\ref{ch11:prop:closed-complete} it is complete. A complete space need not have a right-hand end; what matters is that no limit is missing.}

\section{A geometric bound controls an entire tail}\label{ch11:sec:geometric}
The harmonic sums of Section~\ref{ch11:sec:cauchy} failed to be Cauchy because their steps, though small, added up to a lot. This section shows that steps which shrink by a fixed factor at every stage cannot add up to much. The tool is the finite geometric sum.

\subsection*{The finite geometric sum}
Let $q$ be a real number with $q\ne1$, and let $r$ be a positive integer. Then
\[
1+q+q^2+\cdots+q^{r-1}=\frac{1-q^r}{1-q}.
\]
To see why, first take $r=4$ and call the sum $S=1+q+q^2+q^3$. Multiplying by $q$ raises every power by one:
\begin{align*}
 S&=1+q+q^2+q^3,\\
 qS&=\phantom{1+{}}q+q^2+q^3+q^4.
\end{align*}
When the second line is subtracted from the first, the terms $q$, $q^2$ and $q^3$ cancel, leaving $S-qS=1-q^4$. The left side is $(1-q)S$, and $1-q\ne0$, so we may divide by it: $S=(1-q^4)/(1-q)$. For a general $r$, the same subtraction applies to $S=1+q+\cdots+q^{r-1}$ and $qS=q+q^2+\cdots+q^r$. The terms $q$ through $q^{r-1}$ appear in both, so only $1$ and $q^r$ remain, and $(1-q)S=1-q^r$. As a check, take $q=1/2$ and $r=4$. The sum is $1+\tfrac12+\tfrac14+\tfrac18=\tfrac{15}8$, and the formula gives $\bigl(1-\tfrac1{16}\bigr)\big/\bigl(1-\tfrac12\bigr)=\tfrac{15}{16}\cdot2=\tfrac{15}8$.

From now on let $0<q<1$. We need sums that start at a later power. For integers $m>n\ge0$, factor out $q^n$ and apply the formula with $r=m-n$:
\[
\sum_{j=n}^{m-1}q^j=q^n\bigl(1+q+\cdots+q^{m-n-1}\bigr)=q^n\,\frac{1-q^{m-n}}{1-q}\le\frac{q^n}{1-q}.
\]
The final inequality holds because $q^{m-n}>0$, so $1-q^{m-n}<1$, and the factor $q^n/(1-q)$ is positive. The bound on the right does not depend on $m$: however many terms we add, a sum of consecutive powers starting at $q^n$ never exceeds $q^n/(1-q)$. We call this the \emph{geometric tail bound}.

\subsection*{Powers of a number between zero and one tend to zero}
The bound $q^n/(1-q)$ is useful only if it becomes small as $n$ grows, so we need $q^n\to0$. Section~\ref{ch11:sec:convergence} handled $q=1/2$ by proving $2^n\ge n+1$. For a general $q$ with $0<q<1$, we imitate that proof with $1/q$ in place of $2$. Put $c=1/q-1$. Since $0<q<1$, we have $1/q>1$, so $c>0$, and $1/q=1+c$. We prove by induction that
\[
(1+c)^n\ge1+nc\qquad\text{for every }n\in\NN.
\]
This is known as \emph{Bernoulli's inequality}. For $c=1$ it is the inequality $2^n\ge n+1$ of Section~\ref{ch11:sec:convergence}. At $n=0$ both sides equal $1$. If $(1+c)^n\ge1+nc$ for some $n$, multiply both sides by the positive number $1+c$:
\[
 (1+c)^{n+1}\ge(1+nc)(1+c)
 =1+c+nc+nc^2=1+(n+1)c+nc^2\ge1+(n+1)c,
\]
because $nc^2\ge0$. This is the inequality for $n+1$, so it holds for every $n$. Since $q^n=1/(1+c)^n$, and both $(1+c)^n$ and $1+nc$ are positive, taking reciprocals reverses the order and gives
\[
 0<q^n\le\frac1{1+nc}.
\]
Now let $\eps>0$ be given, and use the Archimedean property to choose a positive integer $N>1/(c\eps)$. If $n\ge N$, then $1+nc>nc\ge Nc>1/\eps$, where the last step multiplies $N>1/(c\eps)$ by the positive number $c$. Taking reciprocals once more gives $q^n\le1/(1+nc)<\eps$. This proves $q^n\to0$.

The threshold depends on $\eps$ and, through $c$, on $q$: the closer $q$ is to $1$, the smaller $c$ is, and the larger the number $1/(c\eps)$ that $N$ must exceed. For example, $q=0.9$ gives $c=1/9$, and for $\eps=0.01$ the recipe asks for $N>900$. In fact $0.9^n<0.01$ already holds from $n=44$ on, since $0.9^{44}\approx0.0097$ while $0.9^{43}\approx0.0108$. As in Section~\ref{ch11:sec:convergence}, a threshold produced by a proof may be far from the smallest one, and that is allowed.

\subsection*{Geometrically shrinking steps give a Cauchy sequence}
\par\noindent\begin{minipage}{\linewidth}
\begin{proposition}\label{ch11:prop:geometric-cauchy}
Let $x_0,x_1,x_2,\ldots$ be a sequence in a metric space $(X,d)$, and suppose there are numbers $D\ge0$ and $0<q<1$ such that
\[
d(x_{j+1},x_j)\le Dq^j\qquad\text{for every }j\in\NN.
\]
Then $d(x_m,x_n)\le Dq^n/(1-q)$ whenever $m>n$, and the sequence is a Cauchy sequence.
\end{proposition}
\end{minipage}\par

Before the proof, here is an example. For the sequence $x_n=1-2^{-n}$ of Section~\ref{ch11:sec:convergence}, each step is
\[
x_{j+1}-x_j=\bigl(1-2^{-(j+1)}\bigr)-\bigl(1-2^{-j}\bigr)=2^{-j}-2^{-(j+1)}=2^{-(j+1)}=\tfrac12\bigl(\tfrac12\bigr)^j.
\]
This number is positive, so it equals its absolute value $d(x_{j+1},x_j)$, and the hypothesis holds with $D=1/2$ and $q=1/2$. The proposition then promises $|x_m-x_n|\le\tfrac12\bigl(\tfrac12\bigr)^n\big/\tfrac12=2^{-n}$ for all $m>n$. Directly, $|x_m-x_n|=2^{-n}-2^{-m}$, which is indeed less than $2^{-n}$.

\begin{proof}
\emph{Step 1: from single steps to many steps.} Fix $n$. We show by induction on $m$ that, for every $m>n$,
\[
d(x_m,x_n)\le\sum_{j=n}^{m-1}d(x_{j+1},x_j).
\]
In words, the distance from $x_n$ to $x_m$ is at most the total length of the single steps between them. For instance, with $m=n+3$ the claim reads
\[
d(x_{n+3},x_n)\le d(x_{n+3},x_{n+2})+d(x_{n+2},x_{n+1})+d(x_{n+1},x_n).
\]
When $m=n+1$, the sum has just one term, $d(x_{n+1},x_n)$, and the claim is an equality. Suppose the claim holds for some $m>n$. The triangle inequality with $x_m$ as the stopping point gives
\[
d(x_{m+1},x_n)\le d(x_{m+1},x_m)+d(x_m,x_n)\le d(x_{m+1},x_m)+\sum_{j=n}^{m-1}d(x_{j+1},x_j)=\sum_{j=n}^{m}d(x_{j+1},x_j),
\]
which is the claim for $m+1$.

\emph{Step 2: a bound for the whole tail.} Combine Step~1 with the assumed bounds on the steps and with the geometric tail bound. For $m>n$,
\[
d(x_m,x_n)\le\sum_{j=n}^{m-1}d(x_{j+1},x_j)\le\sum_{j=n}^{m-1}Dq^j=D\sum_{j=n}^{m-1}q^j\le\frac{Dq^n}{1-q}.
\]
The last step multiplies the geometric tail bound $\sum_{j=n}^{m-1}q^j\le q^n/(1-q)$ by the number $D$, which keeps the inequality in the same direction because $D\ge0$.

\emph{Step 3: the Cauchy condition.} Let $\eps>0$. If $D=0$, Step~2 gives $d(x_m,x_n)\le0$, and hence $d(x_m,x_n)=0$, for all $m>n$, because distances are never negative. So all the terms are equal, and the threshold $N=0$ works. Now let $D>0$. The number $\eps(1-q)/D$ is positive, and $q^n\to0$, so there is a threshold $N$ such that $q^n<\eps(1-q)/D$ for every $n\ge N$. Multiplying by the positive number $D/(1-q)$ gives $Dq^n/(1-q)<\eps$ for every $n\ge N$. Take any two indices $m,n\ge N$. If $m=n$, their distance is $0$, which is less than $\eps$. Otherwise, by symmetry of $d$, we may name them so that $m>n$, and Step~2 gives $d(x_m,x_n)\le Dq^n/(1-q)<\eps$. So the sequence is Cauchy.
\end{proof}

The bound $Dq^n/(1-q)$ in Step~2 involves $n$ but not $m$, so a single choice of $n$ controls the whole tail from $x_n$ onward, however far beyond $n$ the index $m$ goes. The harmonic sums had no such control: there the steps from index $n$ to index $2n$ add up to at least $1/2$, however large $n$ is, while here the steps after index $n$ add up to at most $Dq^n/(1-q)$, which is as small as we like once $n$ is large.

If the space is complete, as the real line and every closed interval are, then the definition of completeness turns the conclusion of the proposition into convergence: the sequence has a limit in the space. This combination of a geometric bound on the steps with completeness is the heart of the contraction theorem in Chapter~\ref{ch:contraction}.

\section*{Exercises}
Exercises~\ref{ex:11.1}--\ref{ex:11.4} are core exercises, and Exercises~\ref{ex:11.5} and~\ref{ex:11.6} are transfer exercises. Complete solutions are in Appendix~\ref{app:solb}, under ``Week~11.''
\begin{exercise}\label{ex:11.1}Suppose $|x-a|\le0.02$ and $|y-b|\le0.03$. Find an upper bound for $|(2x-y)-(2a-b)|$ that holds for all such numbers, and give values of $x$, $y$, $a$ and $b$ for which your bound is attained.\end{exercise}
\begin{exercise}\label{ex:11.2}Prove that the sequence $a_n=3/(n+1)$, for $n\in\NN$, converges to $0$. For each tolerance $\eps>0$, state explicitly the threshold $N$ that you use.\end{exercise}
\begin{exercise}\label{ex:11.3}Show that the function $d_\infty$ on $\RR^n$ from Section~\ref{ch11:sec:metric} satisfies the first two requirements of Definition~\ref{def:metric}: $d_\infty(x,y)\ge0$, with equality exactly when $x=y$, and $d_\infty(x,y)=d_\infty(y,x)$. Together with the triangle inequality proved in that section, this shows that $d_\infty$ is a metric.\end{exercise}
\begin{exercise}\label{ex:11.4}A sequence $x_0,x_1,x_2,\ldots$ in a metric space satisfies $d(x_{j+1},x_j)\le2(1/3)^j$ for every $j\in\NN$. Find an upper bound for $d(x_m,x_n)$ that holds for all $m>n$ and does not depend on $m$. Then show that the sequence is Cauchy, with an explicit threshold for each tolerance $\eps>0$.\end{exercise}
\begin{exercise}\label{ex:11.5}Give a Cauchy sequence in the interval $(2,3)$, with distance $|x-y|$, that has no limit in $(2,3)$. Prove both properties of your sequence.\end{exercise}
\begin{exercise}\label{ex:11.6}Let $A$ be a set of real numbers. A real number $b$ is an \emph{upper bound} for $A$ if $a\le b$ for every $a\in A$. A \emph{maximum} of $A$ is an upper bound for $A$ that is itself a member of $A$. A \emph{supremum}, or \emph{least upper bound}, of $A$ is an upper bound $s$ for $A$ such that $s\le b$ for every upper bound $b$ for $A$.
\begin{enumerate}
\item Show that $1$ is a supremum of the interval $(0,1)$.
\item Show that $(0,1)$ has no maximum.
\item The completeness of the real line is often stated in the following form: \emph{every nonempty set of real numbers that has an upper bound has a supremum}. Explain why this statement would become false if the word ``supremum'' were replaced by ``maximum.''
\end{enumerate}
\end{exercise}

